%==============================================================================
\section{Khái niệm nền --- Trực quan hoá dữ liệu trong ML}
\label{sec:khai-niem-truc-quan}
%==============================================================================

\begin{dinhnghia}[title={Trực quan hoá dữ liệu (Data Visualization)}]
Trực quan hoá dữ liệu là biểu diễn đồ họa của dữ liệu và thông tin.
Nhờ đó ta thấy dữ liệu ``trông như thế nào'' và mức độ tương quan giữa các thuộc tính --- đây là cách nhanh nhất để kiểm tra feature có tương ứng với đầu ra hay không.
\end{dinhnghia}

\begin{ynghia}[title={Vai trò trong pipeline ML}]
\begin{itemize}
  \item \textbf{Khám phá dữ liệu}: pattern, tương quan, outlier; phát hiện missing, sai kiểu, không nhất quán.
  \item \textbf{Chọn đặc trưng}: xác định biến mạnh/yếu so với biến mục tiêu.
  \item \textbf{Đánh giá mô hình}: ROC, precision--recall, ma trận nhầm lẫn.
  \item \textbf{Truyền đạt}: biểu đồ giúp stakeholder không chuyên hiểu kết quả.
\end{itemize}
\end{ynghia}

\begin{phuongphap}[title={Hai nhóm biểu đồ trong phần này}]
\begin{itemize}
  \item \textbf{Univariate} (một biến): histogram, density plot, boxplot --- xem từng thuộc tính độc lập.
  \item \textbf{Multivariate} (nhiều biến): ma trận tương quan, scatter matrix --- xem tương tác giữa các thuộc tính.
\end{itemize}
\end{phuongphap}

\begin{phuongphap}[title={Thư viện Python thường dùng}]
Matplotlib, Seaborn, Plotly, Bokeh.
\end{phuongphap}

\begin{vidu}[title={Tổng quan kỹ thuật trực quan}]
\tsimg{01_data_visualization/img01.jpg}
\end{vidu}
